\documentclass[11pt]{article}

\usepackage[review]{acl}
\usepackage{times}
\usepackage{latexsym}
\usepackage[T1]{fontenc}
\usepackage[utf8]{inputenc}
\usepackage{microtype}
\usepackage{inconsolata}
\usepackage{graphicx}
\usepackage{booktabs}
\usepackage{multirow}
\usepackage{amsmath}

\title{Czech Clinical Named Entity Recognition for MultiClinAI:\\
Morphological Augmentation, CRF Decoding, and Ensemble Voting}

\author{Anonymous}

\begin{document}
\maketitle

\begin{abstract}
We present our system for Czech clinical Named Entity Recognition (NER) submitted to the MultiClinAI
shared task, targeting three entity types: disease, symptom, and procedure. Czech poses particular challenges
as a morphologically rich, inflectional language with limited annotated clinical resources. Our approach
fine-tunes multilingual (XLM-RoBERTa-base) and Czech-specific (RobeCzech) transformer models, optionally
combined with a Conditional Random Field (CRF) decoding layer. The most impactful contribution is a
morphological synonym augmentation strategy that generates additional training documents by replacing rare
entity surface forms with morphological variants from a curated Czech medical synonym dictionary, yielding
a consistent +4\% F1 improvement across all entity types. We further improve performance through weighted
ensemble voting across multiple model configurations and identify a tokenizer artifact in XLM-RoBERTa's
SentencePiece byte-fallback that produced spurious short entity spans, whose resolution brought a substantial
gain on the test set. Our best system achieves strict F1 scores of 0.691 (disease), 0.662 (procedure), and
0.634 (symptom) on the official MultiClinNER Czech test set.
\end{abstract}

% -----------------------------------------------------------------------
\section{Introduction}
% -----------------------------------------------------------------------

Clinical documents encode a wealth of medical knowledge, but extracting structured information from
free-text narratives requires dedicated natural language processing (NLP) systems. Named Entity Recognition
(NER) is a foundational step in clinical information extraction, enabling downstream tasks such as diagnosis
coding, pharmacovigilance, and population health analysis. The MultiClinAI shared task~\cite{multiclinai2026}
challenges participants to build NER systems for seven languages, extracting three entity types
(\textsc{Disease}, \textsc{Symptom}, \textsc{Procedure}) from clinical case reports.

Czech presents distinct challenges compared to other task languages. As a highly inflectional Slavic
language, the same medical concept may appear in many morphological forms depending on grammatical
case, number, and gender. For example, the disease \textit{hypertenze} (hypertension) appears as
\textit{hypertenz\'{i}} (genitive), \textit{hypertenzi} (dative/accusative), or \textit{hypertenz\'{i}}
(instrumental) in natural text. Standard surface-form matching and even pretrained language models
may not generalise well across all inflected forms if they are underrepresented in training data.
Furthermore, the MultiClinNER Czech training set comprises approximately 1,006 documents per entity type,
which is relatively small for fine-tuning large transformer models reliably.

This paper describes our system for the Czech track of MultiClinNER (Subtask 1: entity detection).
Our main contributions are:

\begin{itemize}
  \item A \textbf{morphological synonym augmentation} strategy targeting rare entity surface forms
        that generates 2,161--2,271 additional training documents per entity type, improving
        F1 by approximately 4 percentage points.
  \item A systematic comparison of multilingual (XLM-RoBERTa-base) vs.\ Czech-specific
        (RobeCzech) transformer models, with and without CRF decoding.
  \item A \textbf{sentence-level training} procedure with Czech-aware abbreviation detection
        that improves training efficiency while maintaining entity/context balance.
  \item Identification and correction of a \textbf{SentencePiece tokenizer artifact} that
        caused contradictory BIO label sequences and spurious short-span predictions.
  \item A \textbf{weighted ensemble voting} strategy combining fine-tuned models with a
        GPT-4.1-mini zero-shot model for complementary recall.
\end{itemize}

% -----------------------------------------------------------------------
\section{Related Work}
% -----------------------------------------------------------------------

\paragraph{Transformer-based NER.}
Transformer models pretrained on large corpora~\cite{devlin-etal-2019-bert} have become the
standard for sequence labeling tasks. Token classification with a linear head over contextual
embeddings achieves strong results on many NER benchmarks. Multilingual models such as
mBERT and XLM-RoBERTa~\cite{conneau-etal-2020-unsupervised} enable transfer across languages,
which is especially valuable in low-resource settings.

\paragraph{CRF decoding.}
Conditional Random Fields (CRFs)~\cite{lafferty2001conditional} model label transition
dependencies explicitly, enforcing globally consistent BIO sequences via Viterbi decoding.
Combined with transformer encoders~\cite{ma-hovy-2016-end}, CRF heads prevent invalid
label sequences (e.g., \texttt{I-X} following \texttt{O}) and can improve precision in
span-boundary detection.

\paragraph{Clinical and biomedical NER.}
Clinical NER has benefited substantially from domain-adaptive pretraining. Models
pretrained on biomedical text such as BioBERT~\cite{lee2020biobert} and
clinical corpora outperform general-domain models on tasks like disease and drug extraction.
For Spanish clinical NER, CLIN-X-ES~\cite{gutierrez-fandiño-etal-2022-maria} demonstrated
the value of language-and-domain-specific pretraining. The Czech language has seen less work
in the clinical domain; RobeCzech~\cite{straka-strakova-2020-robeczech} is the primary
Czech monolingual language model, pretrained on large web and news corpora.

\paragraph{Data augmentation for NER.}
The scarcity of annotated clinical data motivates augmentation approaches. \citet{wei-zou-2019-eda}
propose simple operations (synonym replacement, insertion, deletion, swap) that improve text
classification. For NER specifically, synonym substitution must preserve entity boundaries and
annotation offsets. \citet{dai-adel-2020-analysis} analyse various NER augmentation
strategies and find that synonym replacement is among the most effective. We extend this
approach with morphologically varied synonyms appropriate for an inflectional language.

\paragraph{LLM-based NER.}
Large Language Models (LLMs) applied zero-shot or few-shot to NER have shown competitive
performance in general domains~\cite{wang-etal-2023-gpt}, but tend to underperform supervised
models on specialised clinical terminology. Their high precision on rare entities makes them
useful as diversity sources in ensemble systems.

% -----------------------------------------------------------------------
\section{Data}
% -----------------------------------------------------------------------

\subsection{MultiClinNER Czech Dataset}

The MultiClinNER Czech training data is derived from the DisTEMIST, SympTEMIST, and MedProcNER
corpora translated and annotation-projected to Czech~\cite{multiclinai2026}. Documents are
clinical case reports annotated in BRAT standoff format with character-offset entity spans.
Each entity type (disease, symptom, procedure) has a separate set of documents; per the
task guidelines, texts are not shared across entity types due to translation revision
differences, precluding multi-label training.

Table~\ref{tab:data_stats} summarises the dataset statistics for Czech.

\begin{table}[h]
\centering
\small
\begin{tabular}{lrrrr}
\toprule
\textbf{Entity} & \multicolumn{2}{c}{\textbf{Documents}} & \multicolumn{2}{c}{\textbf{Entities}} \\
\cmidrule(lr){2-3}\cmidrule(lr){4-5}
 & \textbf{Train} & \textbf{Dev} & \textbf{Train} & \textbf{Dev} \\
\midrule
Disease   & 1,006 & 252 & 21,083 & 4,710 \\
Symptom   & 1,006 & 252 & 22,525 & 5,281 \\
Procedure & 1,006 & 252 & 22,438 & 5,063 \\
\bottomrule
\end{tabular}
\caption{MultiClinNER Czech dataset statistics (entity counts from gold \texttt{.ann} files).}
\label{tab:data_stats}
\end{table}

\subsection{Morphological Synonym Augmentation}

To address the morphological richness of Czech and the limited training set size, we
constructed a Czech medical synonym dictionary mapping entity surface forms to
morphological variants. The dictionary was built by combining:
(1) manually curated morphological paradigms for Czech medical terms,
and (2) Wikidata concept labels for Czech medical entities.
The final dictionary contains 1,399 synonym entries spanning 887 distinct morphological
families.

Augmented documents are generated as follows. For each entity surface form in the training
set that occurs at most 5 times (rare entities), we look up all available synonyms and
generate one new document per synonym by replacing the entity with the synonym and
adjusting all downstream character offsets accordingly. Entities with
more than 5 occurrences are considered well-represented and excluded to avoid over-amplifying
already common forms. This process yields 2,161 new disease documents, 2,271 new symptom
documents, and 2,271 new procedure documents, more than tripling the training set size.

\subsection{LLM Paraphrase Augmentation}

We additionally generated paraphrases using GPT-4.1-mini following the BioSynNER
methodology, where the model is prompted to rewrite clinical sentences while preserving
entity spans inline. This produced supplementary paraphrase data at approximately \$3.30
per entity type. The paraphrase augmentation provided a small additional improvement
(0.002--0.004 F1 on development set) beyond morphological augmentation.

% -----------------------------------------------------------------------
\section{Methods}
% -----------------------------------------------------------------------

\subsection{Base Models}

We evaluate two pretrained transformer models for Czech:

\textbf{XLM-RoBERTa-base}~\cite{conneau-etal-2020-unsupervised} (\texttt{FacebookAI/xlm-roberta-base}):
A multilingual model pretrained on 2.5TB of filtered CommonCrawl data across 100 languages
using masked language modelling. Despite not being specialised for Czech or clinical text,
its large multilingual pretraining vocabulary handles Czech morphological forms well through
SentencePiece subword tokenisation.

\textbf{RobeCzech}~\cite{straka-strakova-2020-robeczech} (\texttt{ufal/robeczech-base}):
A Czech-specific RoBERTa model pretrained on Czech web text, news, and Wikipedia. Its
Czech-optimised vocabulary produces fewer subword splits per token for Czech words,
potentially yielding better token-level entity alignment.

\subsection{Model Architecture}

Both models are fine-tuned for token classification in a BIO (Begin-Inside-Outside) tagging
scheme, with one model trained per entity type. Separate models for each entity type are
used because the task provides separate annotated text versions per entity.

The architecture consists of:
\begin{enumerate}
  \item The pretrained transformer encoder producing contextual token representations.
  \item A dropout layer (rate 0.1) for regularisation.
  \item A linear projection layer mapping hidden states to the label space (7 labels:
        \texttt{B-TYPE}, \texttt{I-TYPE}, \texttt{O} for each entity type, plus padding label).
  \item Optionally, a \textbf{Conditional Random Field (CRF)} layer replacing the
        standard argmax decoder with Viterbi-decoded label sequences.
\end{enumerate}

During training with CRF, we minimise the negative log-likelihood of the gold label
sequence. At inference, Viterbi decoding finds the globally optimal BIO sequence,
preventing invalid transitions (e.g., \texttt{I-X} after \texttt{O}).

Special tokens (\texttt{[CLS]}, \texttt{[SEP]}) and non-first subword tokens receive
label $-100$ and are excluded from both CRF and cross-entropy loss computation.
Only the first subword token of each word receives a valid label; subsequent subword
tokens inherit the outside label for loss purposes.

\subsection{Tokenizer Artifact Fix}

During development, we identified a critical artifact in XLM-RoBERTa's SentencePiece
byte-fallback tokenisation: for certain byte sequences (typically non-ASCII characters in
rare contexts), the tokeniser emits two consecutive tokens at identical character start
offsets. This occurred approximately 4 times per document in the Czech clinical corpus.

When the standard label alignment algorithm processes these overlapping tokens, it assigns
a \texttt{B-TYPE} label to both tokens at the same position, creating a
\texttt{B}$\to$\texttt{B} transition in the label sequence that contradicts BIO semantics.
During training, this introduces inconsistent supervision; during inference, it causes the
entity reconstruction to close one span and immediately open another, often producing 1--2
character spans corresponding to subword fragments rather than meaningful clinical terms.

We resolve this by:
\begin{itemize}
  \item \textbf{Training}: Overlapping tokens (same start offset as the preceding token)
        receive label $-100$ and are excluded from the loss.
  \item \textbf{Inference}: When an overlapping token is encountered, the entity span
        is extended to the end of the overlapping token rather than split.
  \item \textbf{Post-filtering}: All predicted spans of 1 character are discarded.
        Spans of exactly 2 characters are retained only if they are fully uppercase
        (valid medical abbreviations such as \textit{DM}, \textit{HT}).
\end{itemize}

Table~\ref{tab:bugfix} shows the impact of this fix on the test set.

\subsection{Training Setup}

We train all models with the following hyperparameters: learning rate $2 \times 10^{-5}$
(AdamW), weight decay 0.01, linear warmup for 10\% of steps, batch size 16, maximum
sequence length 512 tokens, and early stopping with patience 3 (monitoring development
F1). Final submission models are retrained on the combined train+dev split for a fixed
10 epochs (no early stopping).

\paragraph{Sentence-level training.}
Rather than processing full documents (which may have variable lengths), we split each
document into sentences and train on sentence-level examples. This increases the
effective batch diversity and entity-to-token ratio per batch. Sentence boundaries are
detected using punctuation followed by whitespace and an uppercase letter, with a
curated list of Czech abbreviations (e.g., \textit{mudr.}, \textit{např.}, \textit{tj.})
that suppress false splits, and a decimal-number check to avoid splitting on numerical
values. We retain pure-\texttt{O} sentences in training to prevent the model from
developing a recall bias (empirically, removing them drops precision to 0.418).

\paragraph{Oversampling.}
Documents containing at least one entity annotation are oversampled once (1$\times$) in
the training set before merging with augmented data. This further increases the proportion
of entity-containing examples without over-weighting augmented data.

\subsection{Ensemble Method}

We combine multiple independently trained model checkpoints through span-level voting.
Each model predicts a set of (start, end, type) entity spans; the ensemble aggregates
these predictions using one of several strategies explored in a development-set sweep:

\begin{itemize}
  \item \textbf{Threshold-$K$}: Accept a span if at least $K$ models predict it.
  \item \textbf{Weighted voting}: Each model is assigned a weight proportional to its
        development F1 score. A span is accepted if the sum of voting model weights
        exceeds a threshold $t$ (swept over $[0.3, 0.8]$).
  \item \textbf{Two-tier}: A core agreement layer (at least $K$ models) plus an
        additional trust layer accepting unique predictions from high-precision models.
\end{itemize}

For disease, the best development-set ensemble combines four models: XLM-R+morph+oversample,
RobeCzech+morph+paraphrase, RobeCzech+morph, and GPT-4.1-mini (zero-shot). Despite
GPT-4.1-mini's low individual F1 (0.384), its predictions have 65.8\% complementarity
with the fine-tuned models: it finds 224 entities missed by all fine-tuned models,
including rare terms, complex multi-word spans, and abbreviations. This boosts ensemble
recall from 0.719 to 0.730 with weighted voting at $t=0.4$.

% -----------------------------------------------------------------------
\section{Experiments and Results}
% -----------------------------------------------------------------------

\subsection{Ablation Study (Development Set)}

Table~\ref{tab:ablation} presents the development-set ablation for disease NER, comparing
the effect of model choice, CRF decoding, and augmentation strategy. Results for
symptom and procedure show consistent trends.

\begin{table}[h]
\centering
\small
\begin{tabular}{llllr}
\toprule
\textbf{Model} & \textbf{Aug.} & \textbf{Head} & \textbf{Extra} & \textbf{F1} \\
\midrule
XLM-R-base & None    & Softmax & {--}         & 0.657 \\
XLM-R-base & Morph   & Softmax & {--}         & 0.710 \\
XLM-R-base & Morph   & CRF     & {--}         & 0.712 \\
XLM-R-base & Morph   & Softmax & OS$\times$1 & 0.710 \\
XLM-R-base & Morph+Para & Softmax & {--}      & 0.704 \\
RobeCzech  & Morph   & Softmax & {--}         & 0.698 \\
RobeCzech  & Morph+Para & Softmax & {--}      & \textbf{0.714} \\
\bottomrule
\end{tabular}
\caption{Development set disease F1 ablation. Aug.\ = augmentation (Morph = morphological
synonyms, Para = LLM paraphrase); Extra: OS$\times$1 = 1$\times$ entity-doc oversampling.
XLM-R-base = \texttt{FacebookAI/xlm-roberta-base}; RobeCzech = \texttt{ufal/robeczech-base}.}
\label{tab:ablation}
\end{table}

Key findings from the ablation:

\textbf{Morphological augmentation} is the single most effective intervention, adding
approximately 4 percentage points F1 over the non-augmented baseline (0.657 $\to$ 0.710--0.712).

\textbf{CRF decoding} provides a small but consistent improvement over softmax for the
same model and augmentation (XLM-R + morph: 0.710 $\to$ 0.712), by enforcing valid
BIO transitions during decoding.

\textbf{RobeCzech vs.\ XLM-RoBERTa}: Despite being Czech-specific, RobeCzech
slightly underperforms XLM-RoBERTa with morphological augmentation alone (0.698 vs.\ 0.710).
However, RobeCzech benefits more from LLM paraphrase augmentation, reaching
0.714 to become the best single model---suggesting its Czech-optimised tokenisation
extracts more value from paraphrased clinical text.

\textbf{LLM paraphrases} add a marginal +0.004 improvement beyond morphological augmentation
alone, indicating diminishing returns when the morphological variants already cover
the most impactful diversity.

\textbf{Oversampling} entity-containing documents (OS$\times$1) yields no additional gain
beyond morphological augmentation, though it was beneficial without augmentation.

\subsection{Ensemble Results (Development Set)}

\begin{table}[h]
\centering
\small
\begin{tabular}{llrrr}
\toprule
\textbf{System} & \textbf{Strategy} & \textbf{P} & \textbf{R} & \textbf{F1} \\
\midrule
\multirow{4}{*}{4-model}
  & Weighted $t$=0.4  & 0.747 & 0.730 & \textbf{0.738} \\
  & Threshold $K$=2   & 0.743 & 0.729 & 0.736 \\
  & Majority $K$=3    & 0.820 & 0.629 & 0.712 \\
  & Union $K$=1       & 0.694 & 0.731 & 0.712 \\
\midrule
3-model & Threshold $K$=2  & 0.753 & 0.719 & 0.736 \\
\midrule
Best single      & {--}             & 0.718 & 0.711 & 0.714 \\
\bottomrule
\end{tabular}
\caption{Disease development set ensemble results (dev set, 4,710 gold entities).
4-model: XLM-R+morph+OS, RobeCzech+morph+para, RobeCzech+morph, GPT-4.1-mini.
3-model: same without GPT-4.1-mini. Best single: RobeCzech+morph+para.}
\label{tab:ensemble}
\end{table}

The weighted voting ensemble with $t$=0.4 achieves F1=0.738, outperforming the best
single model by 2.4 percentage points. Notably, including the GPT-4.1-mini zero-shot
model increases recall (0.719 $\to$ 0.730) without substantially hurting precision (0.753 $\to$ 0.747),
since most of its unique predictions are genuine entities missed by the fine-tuned models.
The 3-model variant (excluding GPT) achieves the same F1 (0.736) through higher precision
(0.753 vs.\ 0.747), providing a useful precision-recall trade-off.

For procedure and symptom, two-model ensembles (XLM-R+CRF and RobeCzech+OS$\times$1)
were used since fewer model variants were trained for these entity types.

\subsection{Bug Fix Impact}

\begin{table}[h]
\centering
\small
\begin{tabular}{llrrr}
\toprule
\textbf{Entity} & \textbf{System} & \textbf{P} & \textbf{R} & \textbf{F1} \\
\midrule
\multirow{2}{*}{Disease}
  & Ensemble (pre-fix) & 0.565 & 0.674 & 0.615 \\
  & Ensemble (post-fix) & 0.708 & 0.675 & \textbf{0.691} \\
\midrule
\multirow{2}{*}{Procedure}
  & Ensemble (pre-fix) & 0.617 & 0.685 & 0.649 \\
  & Ensemble (post-fix) & 0.634 & 0.692 & \textbf{0.662} \\
\midrule
\multirow{2}{*}{Symptom}
  & Ensemble (pre-fix) & 0.564 & 0.629 & 0.595 \\
  & Ensemble (post-fix) & 0.695 & 0.582 & \textbf{0.634} \\
\bottomrule
\end{tabular}
\caption{Test set impact of the SentencePiece tokenizer artifact fix. Pre-fix results from
official evaluation of pre-bugfix submissions; post-fix from bugfix submissions.}
\label{tab:bugfix}
\end{table}

The SentencePiece fix brings substantial improvements, particularly in precision (disease:
0.565 $\to$ 0.708, symptom: 0.564 $\to$ 0.695). This confirms that overlapping token
positions were generating large numbers of spurious short entity spans, inflating false positives.

\subsection{Test Set Results}

\begin{table}[h]
\centering
\small
\begin{tabular}{llrrr}
\toprule
\textbf{Entity} & \textbf{System} & \textbf{P} & \textbf{R} & \textbf{F1} \\
\midrule
\multirow{3}{*}{Disease}
  & Ensemble (4-model) & \textbf{0.708} & \textbf{0.675} & \textbf{0.691} \\
  & XLM-R CRF+morph    & 0.651          & 0.638          & 0.645 \\
  & XLM-R+morph+OS     & 0.650          & 0.639          & 0.644 \\
\midrule
\multirow{2}{*}{Procedure}
  & Ensemble (2-model) & 0.634 & \textbf{0.692} & \textbf{0.662} \\
  & XLM-R CRF+morph    & \textbf{0.668} & 0.655          & 0.662 \\
\midrule
\multirow{2}{*}{Symptom}
  & Ensemble (2-model) & \textbf{0.695} & 0.582 & \textbf{0.634} \\
  & XLM-R CRF+morph    & 0.612          & 0.581 & 0.596 \\
\bottomrule
\end{tabular}
\caption{Official MultiClinNER Czech test set results (strict span matching, post-bugfix).
Ensemble models: Disease = XLM-R×2 + RobeCzech×2; Procedure/Symptom = XLM-R CRF + RobeCzech+OS.}
\label{tab:test}
\end{table}

The ensemble consistently outperforms individual models across all entity types. The
largest gain from ensembling is seen for symptom (+0.038 F1), where the individual CRF
model's lower recall (0.581) is partially recovered by the complementary RobeCzech model.
For procedure, both ensemble and CRF achieve identical F1 (0.662), with the ensemble
recovering higher recall at the cost of lower precision.

Comparing entity types, disease achieves the highest test F1 (0.691), followed by
procedure (0.662) and symptom (0.634). The lower symptom performance likely reflects
the inherently more ambiguous and context-dependent nature of symptom mentions, which
can overlap with disease and procedure terminology.

% -----------------------------------------------------------------------
\section{Error Analysis}
% -----------------------------------------------------------------------

\subsection{GPT-4.1-mini Complementarity}

To understand the sources of ensemble gain from GPT-4.1-mini, we analyse its
predictions on the development set alongside the best fine-tuned model (RobeCzech+morph+para).

\begin{table}[h]
\centering
\small
\begin{tabular}{lcc}
\toprule
\textbf{Metric} & \textbf{RobeCzech} & \textbf{GPT-4.1-mini} \\
\midrule
Individual F1      & 0.714  & 0.384 \\
Entities found     & 3,350 (71.1\%)  & 1,447 (30.7\%) \\
False positives    & 1,318  & 1,431 \\
Unique correct     & 2,127  & \textbf{224} \\
\midrule
Union recall       & \multicolumn{2}{c}{75.9\% (+4.8\% vs.\ single)} \\
Complementarity    & \multicolumn{2}{c}{65.8\%} \\
\bottomrule
\end{tabular}
\caption{Comparison of RobeCzech+morph+para and GPT-4.1-mini on the Czech disease
development set (4,710 gold entities). Percentages are out of 4,710.
Complementarity: share of GPT's correct finds that RobeCzech misses.}
\label{tab:gpt_analysis}
\end{table}

GPT-4.1-mini uniquely recovers 224 disease entities missed by all fine-tuned models.
Analysis of these unique finds reveals three main categories:

\begin{itemize}
  \item \textbf{Rare and complex multi-word terms} (e.g., \textit{mezenchymální tumor},
        \textit{juxtaglomerulární buněčný tumor}): fine-tuned models tend to split or
        partially predict these spans.
  \item \textbf{Medical abbreviations} (e.g., \textit{PKU}): undertrained in the
        limited Czech clinical corpus but known to GPT from broader pretraining.
  \item \textbf{Contextual or implicit disease mentions}: where the model must infer
        disease status from surrounding sentence context.
\end{itemize}

\subsection{Remaining Error Patterns}

Approximately 24\% of gold entities in the disease development set are missed by all
models including GPT-4.1-mini. Manual inspection reveals:

\begin{itemize}
  \item \textbf{Very long entity spans}: Multi-clause noun phrases where span boundaries
        are ambiguous even for annotators.
  \item \textbf{Highly inflected rare forms}: Morphological forms with fewer than 1--2
        training examples even after augmentation.
  \item \textbf{Negated or hypothetical mentions}: Entities appearing in constructions
        such as \textit{bez} (without), \textit{vyloučit} (to rule out), which our
        models treat as positive examples regardless of assertion status.
  \item \textbf{English-Czech code-mixing}: Medical abbreviations written in English
        in Czech clinical text (e.g., \textit{COPD}, \textit{GERD}).
\end{itemize}

False positives are predominantly caused by \textbf{adjacent medical terms} being merged
into a single longer span and \textbf{general medical vocabulary} (anatomical terms, test
names) being incorrectly classified as disease mentions.

% -----------------------------------------------------------------------
\section{Conclusion}
% -----------------------------------------------------------------------

We presented a Czech clinical NER system for the MultiClinAI MultiClinNER shared task,
achieving strict F1 scores of 0.691, 0.662, and 0.634 for disease, procedure, and
symptom respectively. Our main finding is that \textbf{morphological synonym augmentation}
is the most impactful single intervention for Czech clinical NER, providing a consistent
+4\% F1 improvement by generating morphologically varied training documents for rare
entity surface forms. CRF decoding offers small additional gains by enforcing BIO
sequence validity. Ensemble voting across diverse model configurations provides further
gains, with GPT-4.1-mini improving recall on rare and complex entity spans.

We also highlight the importance of \textbf{careful tokenizer handling}: the SentencePiece
byte-fallback artifact in XLM-RoBERTa caused significant precision degradation (precision
dropped to 0.565 for disease) that was recoverable only after diagnosing and fixing
overlapping token alignment. We recommend that practitioners explicitly handle this edge
case when applying XLM-RoBERTa to non-ASCII clinical text.

\textbf{Limitations.} Our system does not address assertion status (negated or hypothetical
entities), which may cause false positives in clinical text. The morphological synonym
dictionary was constructed semi-automatically and may contain errors. Results were
obtained on a single test set and may not generalise across clinical specialities.

% -----------------------------------------------------------------------
\section*{Acknowledgments}
(Removed for review.)

% -----------------------------------------------------------------------
\bibliography{custom}

\end{document}
